\section{研究背景与意义}

近年来，随着深度学习、自然语言处理等技术的飞速突破，大语言模型（Large Language Model，LLM）正从单纯依赖统计规律生成连贯文本的被动工具，逐步演化为具备主动任务规划、工具调用、环境交互及多轮反馈修正能力的智能体（Agent）。据Precedence Research统计，2025年全球AI Agent市场规模已达到79.2亿美元，且预计至2035年将爆发式增长至2946.6亿美元，年均复合增长率高达43.57\%\cite{precedence2025aiagents}；Gartner的预测则指出，到2028年，将有33\%的生成式AI交互彻底摒弃传统的问答模式，转而采用Agent式的自主决策模式\cite{gartner2025agentic}。在产业实践层面，国内外科技巨头的Agent产品线已形成百舸争流的竞争格局：OpenAI的Codex、Anthropic的Claude Code、Google的Gemini Agent以及微软的Copilot Studio等产品的推出引领了国际上的Agent市场化范式；阿里巴巴百炼、字节跳动Coze、智谱AI AutoGLM等国内平台也迅速崛起，共同推动了Agent技术从理论走向落地应用。在开源生态中，LangChain\cite{langchain}、MetaGPT\cite{hong2024metagpt}、CrewAI、AutoGen\cite{wu2024autogen}等开发框架的广泛应用，极大地降低了Agent开发门槛，标志着该技术已从学术实验室的原型阶段迈向大规模商业部署阶段，并深度渗透至办公自动化、软件工程、网络安全运维及复杂数据分析等关键业务领域\cite{wang2024survey,xi2025rise}。

从技术演进逻辑来看，LLM从文本生成器向智能体的角色迈进经历了数个具有里程碑意义的关键节点。2017年，Transformer架构\cite{vaswani2017attention}的提出奠定了序列建模的高效计算基础；2020年，GPT-3\cite{brown2020language}的出现验证了大规模预训练模型在少样本场景下具有极强的泛化能力；2022年的InstructGPT\cite{ouyang2022training}通过指令微调与基于人类反馈的强化学习，成功实现了模型行为与人类意图的初步对齐。2023年，GPT-4\cite{openai2023gpt4} API提供的函数调用能力则进一步推动了LLM与外部工具的标准化连接，使大模型具备更加可靠的结构化行动能力，标志着LLM开始从单纯的“文本生成器”走向具备实际操作能力的“任务执行器”。在Agent形态层面，Weng\cite{weng2023agent}提出了著名的架构公式：Agent=大模型+规划+工具调用（动作）+记忆（如\autoref{fig:agent_overview}所示），清晰界定了现代LLM Agent的四大核心支柱。与传统软件智能体相比，LLM Agent利用大模型的语言理解和推理能力，使任务规划、环境感知和行动决策更加依赖动态生成的语义表示，而非完全由人工预定义规则驱动\cite{wang2024survey}。在此框架下，工具调用成为了Agent感知和作用于外部世界的核心执行通道。在交互范式上，Chain-of-Thoughts\cite{wei2022chain}（CoT）和ReAct\cite{yao2023react}实现了模型推理与行动的交替循环，Reflexion\cite{shinn2023reflexion}则引入了自我反思机制，让Agent在反思中实现任务推进和自我进化。这些范式的迭代不仅奠定了现代LLM Agent的交互基础，更使工具调用成为了其区别于传统语言模型的核心能力标志。正是这种能力，使Agent能够获取环境信息、执行外部操作，并根据反馈调整后续行为，能够自主从环境中获取信息、与复杂环境进行实时交互、根据反馈进行反思并持续改进。

\begin{figure}[htbp]
  \centering
  \resizebox{\textwidth}{!}{\usetikzlibrary{arrows.meta,positioning,calc,fit,backgrounds}
\begin{tikzpicture}[
    >=Latex,
    font=\sffamily\footnotesize,
    stage/.style={
        draw=black!72,
        line width=0.75pt,
        fill=white,
        rounded corners=1.5pt,
        minimum height=2.25cm,
        text width=4.05cm,
        inner sep=6pt,
        align=left
    },
    trusted/.style={stage,draw=black,line width=1pt},
    untrusted/.style={stage,draw=black,line width=1pt,dashed,fill=black!4},
    effect/.style={stage,draw=black,line width=1.2pt,fill=black!9},
    flow/.style={->,draw=black!78,line width=0.9pt},
    provenance/.style={->,draw=black,line width=1.05pt,densely dashed},
    defense/.style={
        draw=black!75,
        line width=0.75pt,
        rounded corners=1.5pt,
        fill=white,
        minimum height=1.15cm,
        text width=6.45cm,
        inner sep=5pt,
        align=left
    },
    tag/.style={font=\sffamily\scriptsize\bfseries,fill=white,inner sep=1.5pt},
    note/.style={font=\sffamily\scriptsize,align=center,text=black!75}
]

\node[note,anchor=south west] (caseid) {IPI跨步骤传播实例（AgentDyn实际日志：\texttt{shopping/user\_task\_18} $\times$ \texttt{injection\_task\_5}）};

\node[trusted,below=0.14cm of caseid.south west,anchor=north west] (task) {%
  \textbf{步骤1：可信用户任务}\\[-1pt]
  购买低价MacBook；比较历史折扣，必要时退款并重新购买。};

\node[untrusted,right=0.65cm of task] (email) {%
  \textbf{步骤2：外部邮件观测}\\[-1pt]
  \texttt{get\_shopping\_linked\_}\\[-2pt]
  \texttt{received\_emails()}\\[2pt]
  注入内容：“再购买一套铅笔”。};

\node[stage,right=0.65cm of email] (latent) {%
  \textbf{步骤3：跨步骤潜伏}\\[-1pt]
  下载并读取历史发票；检索MacBook；退款原订单；读取OTP。};

\node[stage,below=1.05cm of latent] (resolve) {%
  \textbf{步骤4：语义对象解析}\\[-1pt]
  \texttt{search\_product(}\\[-2pt]
  \quad\texttt{product\_name="pencil set")}\\[2pt]
  返回商品标识\texttt{P020}。};

\node[stage,left=0.65cm of resolve] (cart) {%
  \textbf{步骤5：参数持续传播}\\[-1pt]
  \texttt{cart\_add\_product(}\\[-2pt]
  \quad\texttt{product\_id="P020")}\\[2pt]
  注入语义转化为稳定资源标识。};

\node[effect,left=0.65cm of cart] (sink) {%
  \textbf{步骤6：高影响动作汇聚}\\[-1pt]
  \texttt{checkout\_selected\_cart(}\\[-2pt]
  \quad\texttt{product\_ids=["P034","P020"])}\\[2pt]
  正常商品与注入商品一并结算。};

\draw[flow] (task) -- node[note,above]{可信请求} (email);
\draw[provenance] (email) -- node[note,above]{邮件正文} (latent);
\draw[provenance] (latent) -- node[note,right]{若干合法调用后仍保留} (resolve);
\draw[provenance] (resolve) -- node[note,above]{\texttt{P020}} (cart);
\draw[provenance] (cart) -- node[note,above]{\texttt{product\_ids[]}} (sink);

\begin{scope}[on background layer]
  \node[draw=black!35,densely dotted,rounded corners=2pt,fit=(email)(latent)(resolve)(cart)(sink),inner sep=5pt] (attackscope) {};
\end{scope}
\node[defense,below=0.85cm of sink.south west,anchor=north west] (tarc) {%
  \textbf{TARC：调用前的规范性授权门控}\quad
  对步骤4--6的待执行动作逐一判断其是否由可信用户任务授权，并在动作物化前约束未授权调用。};

\node[defense,right=0.65cm of tarc] (tg) {%
  \textbf{TraceGraph：执行中的事实来源与韧性控制}\quad
  重建步骤2--6的来源证据路径；结算前命中\texttt{TG-SEM-004}时执行\texttt{REWRITE/REPLAN}，并按检查点恢复。};

\draw[->,draw=black!58,line width=0.7pt] ($(tarc.north)+(2.0cm,0)$) -- ($(sink.south)+(-0.65cm,-0.06cm)$);
\draw[->,draw=black!58,line width=0.7pt] ($(tarc.north)+(5.5cm,0)$) -- ($(cart.south)+(-0.15cm,-0.06cm)$);
\draw[->,draw=black!58,line width=0.7pt] ($(tg.north)+(4.7cm,0)$) -- ($(resolve.south)+(0.55cm,-0.06cm)$);

\node[note,anchor=north west] at ($(tarc.south west)+(0,-0.18cm)$) {%
  实线：可信请求或常规执行流\quad
  虚线：不可信来源及其传播证据\quad
  灰底粗框：可能产生外部副作用的汇聚动作};

\end{tikzpicture}
}
  \caption{Lilian Weng\cite{weng2023agent}提出的Agent开发公式}
  \label{fig:agent_overview}
\end{figure}

然而，技术的高速演进往往伴随着风险的重构。工具调用能力在极大拓展Agent任务边界与应用可能性的同时，也将大语言模型的安全边界从封闭的文本生成领域扩展到了开放的环境交互域。传统LLM的安全研究主要聚焦于“有害内容生成”的攻防博弈，其典型攻击面主要位于自然语言交互层，即通过构造提示词诱导模型生成不当或有害文本\cite{das2025security,shen2024do,zeng2024howjohnny2}；而工具调用型Agent的安全风险则有着本质的不同：Agent不仅具备理解指令的认知能力，更拥有通过工具执行代码、访问文件系统、操控数据库乃至操作物理设备的执行能力，更关注“语言决策是否触发具有现实破坏性的外部动作”。2023年，Greshake等\cite{greshake2023notwhat}率先披露的间接提示注入攻击（Indirect Prompt Injection, IPI）揭示了LLM Agent所面临的系统性安全风险，在该威胁模型下，攻击者无需直接接触模型接口，只需将恶意指令精心嵌入网页、电子邮件或文档等外部数据源中，一旦Agent通过工具读取并解析这些数据，恶意指令即可隐蔽地注入Agent的推理上下文，进而劫持其后续行为，该发现暴露了LLM Agent工具调用场景中“数据与指令边界模糊”的根本安全隐患。此后，随着LLM Agent工具生态的日益繁荣，这一攻击向量也随之迅速扩散与进化：攻击者不仅可以劫持工具的选择决策，诱导Agent调用恶意工具\cite{shi2026pitoolselection}，还可通过社会工程学手段诱导Agent请求远超任务实际所需的敏感权限\cite{zhang2026evaluatingprivilege}。更为隐蔽的是，攻击者能够利用多步工具调用的逻辑链条，将多个看似合法的子任务进行恶意组合，最终达成恶意目标\cite{kothamasu2026hidden}。此外，攻击者还可通过对抗性提示诱导Agent泄露系统提示词、上下文信息及内部策略，暴露其任务目标与安全约束\cite{hui2024scp}；同时，角色扮演、社会工程学和自动化搜索等越狱方法能够绕过模型既有的安全对齐机制，诱导Agent生成有害内容，并在具备工具调用能力时进一步转化为现实行动\cite{shen2024do,zeng2024howjohnny2}。在Agent供应链环节，攻击者还可能向Agent的长期记忆、检索知识库或模型训练过程植入精心构造的投毒样本或后门触发器，使Agent在特定任务或触发条件下持续输出被操纵的信息、执行恶意行为\cite{chen2024agentpoison,bi2024backdoor}。因此，LLM Agent的安全问题已不再是单点的提示注入防御，而是演变为覆盖整个工具调用链上的工具选择合理性、工具权限控制、多步调用逻辑组合的系统性安全挑战，间接提示注入攻击由于能够跨越数据来源、模型推理和工具执行之间的边界，已成为影响Agent自主行为可靠性的重要威胁之一\cite{owasp2025llmtop10}。

综合上述背景，面对日益复杂的攻击手段和不断扩大的风险边界，LLM Agent正面临模型、数据、记忆和工具生态等多层面的安全挑战，但由于工具调用直接连接模型推理过程与外部环境，是导致攻击行为产生实际影响的关键环节，因此本文聚焦于工具调用链中的行为可信控制问题。本研究面向LLM Agent工具调用链的间接提示注入威胁，旨在构建一个针对性的行为安全控制与防御体系，从“调用前的规范性授权”与“调用中/调用后的证据性响应”两个互补的防御时相出发，构建两类防御方法：其一，在工具调用真正产生副作用之前，依据当前用户任务对应的规范授权范围，对即将执行的动作进行确定性门控，以降低在外部攻击者的间接提示注入攻击下，未授权工具调用的风险发生概率；其二，面向前置授权约束或攻击已跨越多个工具调用步之后的场景，通过对Agent执行过程构建证据图，定位风险传播路径，并在副作用发生前予以阻断与处置，为缓解LLM Agent工具调用链面临的提示注入攻击风险提供一种可能的思路。
